%!TEX program = xelatex
\documentclass[a4paper,14pt]{extarticle}

\usepackage{fontspec}
\setmainfont{Liberation Serif}
\setsansfont{Liberation Sans}
\setmonofont{Liberation Mono}
\newfontfamily\cyrillicfonttt{Liberation Mono}
\usepackage{polyglossia}
\setdefaultlanguage{russian}
\setotherlanguage{english}

\usepackage{geometry}
\geometry{a4paper,left=30mm,right=15mm,top=20mm,bottom=20mm}

\usepackage{setspace}
\onehalfspacing

\usepackage{titlesec}
\titleformat{\section}{\normalfont\large\bfseries}{\thesection.}{0.5em}{}

\usepackage{fancyhdr}
\setlength{\headheight}{17pt}
\pagestyle{fancy}
\fancyhf{}
\fancyhead[R]{\thepage}
\renewcommand{\headrulewidth}{0pt}

\usepackage{graphicx}
\usepackage{caption}
\usepackage{listings}
\usepackage{xcolor}
\usepackage{tabularx}
\usepackage{booktabs}
\usepackage{float}

\lstdefinestyle{PLIstyle}{
  basicstyle=\small\ttfamily,
  numbers=none,
  breaklines=true,
  frame=single,
  showstringspaces=false
}

\usepackage[hidelinks]{hyperref}
\usepackage{enumitem}
\usepackage{amsmath}
\usepackage{indentfirst}
\setlength{\parindent}{1.25cm}

\captionsetup[figure]{name=Рис.,labelsep=period}

\begin{document}

\thispagestyle{empty}
\begin{center}
  \textbf{Санкт-Петербургский политехнический университет}\\
  \textbf{Институт компьютерных наук и технологий}\\
  \textbf{Высшая школа программной инженерии}\\[2cm]

  {\large Отчёт по этапу 2 (модификация синтаксического анализатора)}\\[0.5cm]
  {\large Компилятор с ЯВУ (подмножество PL/I)}\\[1.5cm]

  Дисциплина: «Системы программирования»\\[0.5cm]
  Вариант~№~11\\[1.5cm]
\end{center}

\vfill
\begin{center}
  Санкт-Петербург\\
  2026
\end{center}

\newpage
\tableofcontents
\newpage

\section{Задание и целевая программа}

По согласованию с преподавателем персональный вариант~№~11 задаётся исходной
программой следующего вида (без обязательной разрядности \texttt{(15)} в
объявлениях \texttt{BIN FIXED}; инициализация десятичным литералом \texttt{3}):

\begin{figure}[H]
\centering
\begin{lstlisting}[style=PLIstyle]
EX11: PROC OPTIONS ( MAIN );
DCL A BIN FIXED INIT ( 3 );
DCL B DEC FIXED INIT ( 3 );
DCL C BIN FIXED;
C = B * 5;
C = A = C;
END EX11;
\end{lstlisting}
\caption{Персональный вариант (целевой синтаксис)}
\label{fig:variant}
\end{figure}

Контрольный пример макета --- программа \texttt{examppl.pli} (демонстрация
преподавателя) должна по-прежнему успешно обрабатываться.

\section{Ключевые решения по замечаниям преподавателя}

\subsection{Нетерминалы \texttt{LID}, \texttt{ZNL} и матрица \texttt{TPR}}

В макете столбцы матрицы смежности \texttt{TPR} индексируются только
\textbf{нетерминалами} (число столбцов \texttt{NNETRM}). Символы \texttt{LID}
(десятичный литерал в \texttt{INIT} для \texttt{DEC FIXED}) и \texttt{ZNL}
(логический знак «\texttt{=}» в отличие от арифметических \texttt{ZNK}) должны
иметь номера столбцов \texttt{0..NNETRM-1}. Поэтому в таблице \texttt{VXOD}
нетерминалы \texttt{LID} и \texttt{ZNL} размещаются \textbf{сразу после}
\texttt{ZNK} (индексы 16 и~17), а не в хвосте таблицы: иначе функция
\texttt{numb()} возвращает индекс строки терминала и происходит выход за границы
\texttt{TPR[][0..NNETRM-1]}.

\subsection{Пять вариантов \texttt{ODC}}

Для пункта~8 методички фиксируются пять альтернатив \texttt{<ODC>}:
\texttt{BIN FIXED} без разрядности; \texttt{BIN FIXED} с разрядностью;
\texttt{BIN FIXED} с \texttt{INIT(LIT)} (с разрядностью или без --- в
реализации задаётся ветвями \texttt{SINT}); \texttt{DEC FIXED} без разрядности;
\texttt{DEC FIXED} с разрядностью; объявления \texttt{DEC} с \texttt{INIT(LID)}.
В отчётной грамматике \texttt{INIT} для \texttt{DEC} связывается с
\texttt{LID}, для \texttt{BIN} --- с \texttt{LIT}.

\subsection{Разделение \texttt{ZNK} и \texttt{ZNL}}

Арифметические знаки (\texttt{+}, \texttt{-}, \texttt{*}) остаются в
\texttt{ZNK}. Терминал «\texttt{=}» вводит нетерминал \texttt{ZNL} (правила
\texttt{226--228} в \texttt{SINT[]}) для разбора сравнения в конструкциях вида
\texttt{C = A = C;}.

\subsection{Оператор присваивания и правая часть}

В тексте методички правая часть может обозначаться как \texttt{LVI}; в текущей
реализации макета распознавание выполняется через нетерминал \texttt{AVI} с
расширенными ветвями (в т.\,ч. умножение и вложенное использование
\texttt{=}).

\subsection{Формат входного файла}

Исходный текст читается блоками по 80~байт (\texttt{fread}). Необходимо
представлять исходный текст в формате «карт» фиксированной длины (как в
\texttt{examppl.pli} из \texttt{lab-step-1}), без разрыва лексем на границе
карты.

\section{Модификация грамматики (фрагмент)}

Ниже --- фрагмент правил с выделением изменений (курсивом) по сравнению с базовой
грамматикой шага~1:

\begin{enumerate}
  \item $\langle\text{TEL}\rangle ::= \langle\text{ODC}\rangle \mid \langle\text{TEL}\rangle\langle\text{ODC}\rangle \mid \langle\text{TEL}\rangle\langle\text{OPA}\rangle$
  \item $\langle\text{ODC}\rangle$ --- пять альтернатив: \texttt{BIN FIXED} (с разрядностью или без, с \texttt{INIT(LIT)} или без), \texttt{DEC FIXED} (с разрядностью или без), \texttt{DEC} с \texttt{INIT(LID)}; в реализации задаётся ветвями \texttt{SINT[]}.
  \item $\langle\text{LID}\rangle ::= \langle\text{MAN}\rangle$ (десятичный литерал для \texttt{DEC})
  \item $\langle\text{OPA}\rangle ::= \langle\text{IPE}\rangle\texttt{=}\langle\text{AVI}\rangle\texttt{;}$
\end{enumerate}

\section{Состояние реализации в \texttt{komppl.c}}

В каталоге \texttt{lab-step-2/src/komppl.c} заданы:
\texttt{\#define NSINT 320}, \texttt{NNETRM 18}, \texttt{NVXOD 54}; таблица
\texttt{VXOD} упорядочена так, что первые 18 записей --- нетерминалы
(\texttt{AVI}.. \texttt{ZNK}, \texttt{LID}, \texttt{ZNL}); далее --- терминалы.
Матрица \texttt{TPR} имеет размер \texttt{NV}$\times$\texttt{NNETRM}; в строке
\texttt{TEL} отмечена достижимость \texttt{ZNL}; для термина «\texttt{=}»
указаны и \texttt{ZNK}, и \texttt{ZNL} (по смыслу разных веток разбора).

Тестовые файлы: \texttt{lab-step-2/input/examppl.pli} (копия эталона из
\texttt{lab-step-1}) и \texttt{lab-step-2/input/ex11.pli} (в формате карт
80~символов). Полная целевая строка \texttt{C = A = C;} и объявления без явной
разрядности требуют доводки ветвей \texttt{SINT[]} и семантики и могут
документироваться на следующих этапах.

% Обрезка отчёта: далее по методичке следует листинг с правилами 201--231 в SINT[] --- не включается.

\end{document}
